\documentclass[../../main.tex]{subfiles}

\begin{document}

\section{拓扑空间}\label{section:ch02-拓扑空间}

这里给出的拓扑空间的标准定义是经历了很长时间才形成的。许多数学家，如 Fr\'echet、Hausdorff 等人，在 20 世纪的头十年里给出过不同的定义，但是这些定义很快成为过眼烟云，数学家们很快就确定了最合适的定义。当然，他们希望定义能够有尽可能广泛的包容性，从而把数学中许多有用的例子，如欧氏空间、无穷维欧氏空间及它们之间的函数空间，都作为它的特例包括进去。他们又希望使定义尽可能地狭窄，从而使得那些熟知的空间中的标准定理对于一般拓扑空间也能够成立。当我们试图整理一个新的数学概念时，常常会遇到这样的问题，即如何给出恰如其分的定义。最后确立的拓扑空间的定义看上去有些抽象，但是，当你用不同的方法构造拓扑空间时，将会对它的真正内涵有更好的理解。

\begin{definition}[拓扑与拓扑空间]\label{definition:ch02-拓扑空间}
集合 $X$ 上的一个拓扑（topology）乃是 $X$ 的子集的一个族 $\mathcal{T}$，它满足以下条件：
\begin{enumerate}[(1)]
\item $\varnothing$ 和 $X$ 在 $\mathcal{T}$ 中.
\item $\mathcal{T}$ 的任意子族的元素的并在 $\mathcal{T}$ 中.
\item $\mathcal{T}$ 的任意有限子族的元素的交在 $\mathcal{T}$ 中.
\end{enumerate}
一个指定了拓扑 $\mathcal{T}$ 的集合 $X$ 叫做一个拓扑空间（topological space）。

确切地说，一个拓扑空间就是一个有序偶对 $(X, \mathcal{T})$，其中 $X$ 是一个集合，$\mathcal{T}$ 是 $X$ 上的一个拓扑。在不致混淆的情况下，常常不专门提到 $\mathcal{T}$。

如果 $X$ 是带有拓扑 $\mathcal{T}$ 的一个拓扑空间，我们说 $X$ 的一个子集 $U$ 是 $X$ 中的开集，如果 $U$ 是族 $\mathcal{T}$ 的一个元素。使用这个术语，就可以把拓扑空间说成是一个集合 $X$ 连同它的子集（称之为开集）的一个族，使得 $\varnothing$ 和 $X$ 是开集，开集的任意并和有限交都是开集。
\end{definition}

\par
\begin{figure}[htbp]
\centering
\includegraphics[width=0.55\textwidth]{ch02-fig12-1}
\caption{三元素集上的若干拓扑}\label{figure:ch02-三元素拓扑}
\end{figure}

\par
\begin{figure}[htbp]
\centering
\includegraphics[width=0.5\textwidth]{ch02-fig12-2}
\caption{不是拓扑的两个族}\label{figure:ch02-非拓扑}
\end{figure}

\begin{example}\label{example:ch02-三元素拓扑}
设 $X$ 是三元素集，$X = \{a, b, c\}$。在 $X$ 上有多种可能的拓扑，有几种画在\cref{figure:ch02-三元素拓扑}中，其中右上角图所表示的拓扑中，其开集是 $X$，$\varnothing$，$\{a, b\}$，$\{b\}$ 和 $\{b, c\}$。左上角图所表示的拓扑中只有两个开集 $X$ 和 $\varnothing$。右下角图所表示的拓扑由 $X$ 的所有子集组成。通过 $a, b, c$ 间的置换还可以得到 $X$ 上的另外一些拓扑。


由此可见，即使是只有三个元素的集合，也有许多种不同的拓扑。但这并不是说 $X$ 的子集的任意一个族都是 $X$ 的拓扑。比如，\cref{figure:ch02-非拓扑}所表示的族就不是一个拓扑。

\end{example}

\begin{example}\label{example:ch02-离散与密着拓扑}
如果 $X$ 为任意的一个集合，$X$ 所有子集的族是 $X$ 的一个拓扑，称之为离散拓扑（discrete topology）。仅由 $X$ 和 $\varnothing$ 组成的族也是 $X$ 的一个拓扑，称之为密着拓扑（indiscrete topology）或平庸拓扑（trivial topology）。
\end{example}

\begin{example}\label{example:ch02-有限补拓扑}
设 $X$ 是一个集合，$\mathcal{T}_{f}$ 是使得 $X - U$ 或者是有限集或者等于 $X$ 的那些 $X$ 的子集 $U$ 的全体。那么 $\mathcal{T}_{f}$ 是 $X$ 上的一个拓扑，称之为有限补拓扑（finite complement topology）。因为 $X - X = \varnothing$ 是有限集，$X - \varnothing = X$，所以 $X$ 与 $\varnothing$ 都在 $\mathcal{T}_{f}$ 中。若 $\{U_{\alpha}\}$ 是 $\mathcal{T}_{f}$ 中非空元素的一个加标族，为了证明 $\bigcup U_{\alpha}$ 在 $\mathcal{T}_{f}$ 中，只要确定
\begin{align*}
X - \bigcup U_{\alpha} = \bigcap (X - U_{\alpha}).
\end{align*}
因为每一个集合 $X - U_{\alpha}$ 是有限集，所以右边的集合是一个有限集。如果 $U_{1}, \cdots, U_{n}$ 是 $\mathcal{T}_{f}$ 的非空元素，为了证明 $\bigcap_{i=1}^{n} U_{i}$ 在 $\mathcal{T}_{f}$ 中，只要确定
\begin{align*}
X - \bigcap_{i=1}^{n} U_{i} = \bigcup_{i=1}^{n} (X - U_{i}).
\end{align*}
因为上式右边是有限集的有限并，所以是有限集。
\end{example}

\begin{example}\label{example:ch02-可数补拓扑}
设 $X$ 是一个集合，$\mathcal{T}_{c}$ 是使得 $X - U$ 是可数集或者等于 $X$ 的所有 $X$ 的子集 $U$ 的全体。容易验证，$\mathcal{T}_{c}$ 是 $X$ 上的一个拓扑。
\end{example}

\begin{definition}[细于与粗于]\label{definition:ch02-细于粗于}
设 $\mathcal{T}$ 和 $\mathcal{T}'$ 是给定集合 $X$ 上的两个拓扑，如果 $\mathcal{T}' \supset \mathcal{T}$，则称 $\mathcal{T}'$ 细于（finer）$\mathcal{T}$。若 $\mathcal{T}' \supset \mathcal{T}$ 是真包含关系，则称 $\mathcal{T}'$ 严格细于（strictly finer）$\mathcal{T}$。这两种情形有时也分别称之为 $\mathcal{T}$ 粗于（coarser）$\mathcal{T}'$，和 $\mathcal{T}$ 严格粗于（strictly coarser）$\mathcal{T}'$。我们说 $\mathcal{T}$ 与 $\mathcal{T}'$ 是可比较的，如果或者 $\mathcal{T}' \supset \mathcal{T}$ 或者 $\mathcal{T} \supset \mathcal{T}'$。
\end{definition}

受这个术语的启发，我们不妨把拓扑空间设想成一辆装满石子的卡车，一堆石子以及若干堆石子合在一起都是开集。现在把石子敲成小碎块，那么开集族就变大了。因此，按照这种运算，拓扑，如同石子一样，就变细了。

$X$ 上的两个拓扑不一定可以比较。在前面的\cref{figure:ch02-三元素拓扑}中，右上角拓扑严格细于第一列中三个拓扑中的每一个拓扑，又严格粗于第三列另外两个拓扑中的每一个拓扑。但是它与第二列中的任何一个拓扑都不可比较。

对于这个概念有时也使用其他的术语。有些数学家把 $\mathcal{T}' \supset \mathcal{T}$ 称为 $\mathcal{T}'$ 大于（larger）$\mathcal{T}$，或者 $\mathcal{T}$ 小于（smaller）$\mathcal{T}'$。如果不用那个比较形象的术语“细于”和“粗于”的话，这倒也是一个合适的用语。

有许多数学家对这个概念使用“弱于”和“强于”这种词。不幸的是，他们中有些人（特别是分析学家）往往把 $\mathcal{T}' \supset \mathcal{T}$ 称为 $\mathcal{T}'$ 强于 $\mathcal{T}$，而另一些人（特别是拓扑学家）则称为 $\mathcal{T}'$ 弱于 $\mathcal{T}$。如果读者在有些书中见到“强拓扑”或“弱拓扑”这类的术语，那就必须从文章中去判断它的含义。本书不采用这些术语。

\end{document}